% ============ 第1章 §1.12 Hardy空间(原书 pp.108-114) ============
\section{Hardy 空间}\label{sec:1.12}
\index{Hardy 空间!(R 上的)}
\index{H1 范数}


在 $p = 1$ 的情形,Hilbert 变换在 $L^{1}(\mathbb{R})$ 上不是有界的,但本节引入 Hilbert 变换在其上有界的空间----Hardy 空间 $H^{1}(\mathbb{R})$.设 $f \in L^{1}(\mathbb{R})$,且存在 $g \in L^{1}(\mathbb{R})$ 使 $-i\xi|\xi|^{-1}\hat{f}(\xi) = \hat{g}(\xi)$($\xi \in \mathbb{R}$).设 $\varphi \in \mathcal{S}(\mathbb{R})$,$\int_{\mathbb{R}}\varphi(x)\dd x = 1$,$\varphi_{\varepsilon}(x) = \varphi(x/\varepsilon)/\varepsilon$.由例~\ref{exa:1.1.10} 与定理~\ref{thm:1.1.9},当 $\varepsilon \to 0$ 时,$f * \varphi_{\varepsilon}$,$g * \varphi_{\varepsilon}$ 分别依 $L^{1}(\mathbb{R})$ 范数收敛到 $f$,$g$.另外,因 $f, g, \varphi_{\varepsilon} \in L^{1}(\mathbb{R})$,
\begin{align*}
\mathcal{F}(g * \varphi_{\varepsilon}) & = \mathcal{F}(g)\mathcal{F}(\varphi_{\varepsilon}) = -i\dfrac{\xi}{|\xi|}\mathcal{F}(f)\mathcal{F}(\varphi_{\varepsilon}) \\
& = -i\dfrac{\xi}{|\xi|}\mathcal{F}(f * \varphi_{\varepsilon}) = \mathcal{F}\left(H(f * \varphi_{\varepsilon})\right).
\end{align*}
因 $f, g \in L^{1}(\mathbb{R})$,$\varphi_{\varepsilon} \in L^{2}(\mathbb{R})$,故 $f * \varphi_{\varepsilon}, g * \varphi_{\varepsilon} \in L^{2}(\mathbb{R})$,因此 $H(f * \varphi_{\varepsilon}) \in L^{2}(\mathbb{R})$,
\begin{equation*}
g * \varphi_{\varepsilon} = H(f * \varphi_{\varepsilon}).
\end{equation*}
令 $\varepsilon \to 0$,则左边依 $L^{1}(\mathbb{R})$ 范数收敛到 $g$,右边依测度收敛到 $Hf$.因此 $g = Hf$.

考虑这一点,定义下面的函数空间.

\begin{definition}\label{def:1.12.1}
使 $-i\xi|\xi|^{-1}\hat{f}(\xi) = \hat{g}(\xi)$($\xi \in \mathbb{R}$)成立的 $g \in L^{1}(\mathbb{R})$ 存在的 $f \in L^{1}(\mathbb{R})$ 的全体,称为 Hardy 空间 $H^{1}(\mathbb{R})$.对 $f \in H^{1}(\mathbb{R})$,
\begin{equation*}
\| f \|_{H^{1}(\mathbb{R})} = \| f \|_{H^{1}} = \| f \|_{L^{1}(\mathbb{R})} + \| Hf \|_{L^{1}(\mathbb{R})}
\end{equation*}
称为 $f$ 的 $H^{1}$ 范数.
\end{definition}

$H^{1}$ 范数是范数是显然的.设 $\{f_{j}\}_{j=1}^{\infty}$ 为 $H^{1}(\mathbb{R})$ 的 Cauchy 列.则因 $L^{1}(\mathbb{R})$ 完备,存在 $f, g \in L^{1}(\mathbb{R})$ 使 $\| f_{j} - f \|_{L^{1}} \to 0$,$\| Hf_{j} - g \|_{L^{1}} \to 0$($j \to \infty$).由定理~\ref{thm:1.2.3},
\begin{equation*}
\lim_{j \to \infty}\mathcal{F}f_{j}(\xi) = \mathcal{F}f(\xi), \quad \lim_{j \to \infty}\mathcal{F}(Hf_{j})(\xi) = \mathcal{F}g(\xi).
\end{equation*}
另一方面,因 $\mathcal{F}(Hf_{j})(\xi) = -i\xi|\xi|^{-1}\mathcal{F}f_{j}(\xi)$,
\begin{equation*}
-i\frac{\xi}{|\xi|}\mathcal{F}f(\xi) = \mathcal{F}g(\xi).
\end{equation*}
因此 $f \in H^{1}(\mathbb{R})$ 且 $\lim_{j \to \infty}\| f_{j} - f \|_{H^{1}} = 0$,可知 $H^{1}(\mathbb{R})$ 是 Banach 空间.

若 $f \in L^{1}(\mathbb{R}) \cap L^{2}(\mathbb{R})$ 且 $Hf \in L^{1}(\mathbb{R})$,则因 $f \in L^{2}(\mathbb{R})$,$\mathcal{F}(Hf) = -i\xi|\xi|^{-1}\mathcal{F}(f)$,故 $f \in H^{1}(\mathbb{R})$.下面考虑更具体的例子.

\begin{example}\label{exa:1.12.2}
若 $f \in C^{1}(\mathbb{R})$ 且 $|f'(x)| \leq A_{1}(1 + |x|)^{-2}$,$|f(x)| \leq A_{2}(1 + |x|)^{-3}$($x \in \mathbb{R}$),且 $\int_{\mathbb{R}}f(x)\dd x = 0$,则 $f \in H^{1}(\mathbb{R})$.

因此,特别是当 $f \in \mathcal{S}(\mathbb{R})$ 且 $\int_{\mathbb{R}}f(x)\dd x = 0$ 时,$f \in H^{1}(\mathbb{R})$.

实际上,由引理~\ref{lem:1.3.2},$|Hf(x)| \leq C_{1}(A_{1} + A_{2})(1 + |x|)^{-1}$.当 $|x| \geq 1$ 时,由 $\int_{\mathbb{R}}f(x)\dd x = 0$,
\begin{equation*}
Hf(x) = \frac{1}{\pi}\mathrm{p.v.}\int_{\mathbb{R}}\frac{1}{x - y}f(y)\dd y - \frac{1}{\pi}\int_{\mathbb{R}}\frac{1}{x}f(y)\dd y
\end{equation*}
\begin{equation*}
= \frac{1}{\pi x}\mathrm{p.v.}\int_{\mathbb{R}}\frac{1}{x - y}yf(y)\dd y,
\end{equation*}
故注意 $|(xf(x))'| \leq (A_{1} + A_{2})(1 + |x|)^{-1}$,$|xf(x)| \leq A_{2}(1 + |x|)^{-2}$,再使用引理~\ref{lem:1.3.2},得 $|Hf(x)| \leq 2C_{1}(A_{1} + A_{2})(1 + |x|)^{-2} \in L^{1}(\mathbb{R})$.又 $f \in L^{1}(\mathbb{R}) \cap L^{2}(\mathbb{R})$,故 $f \in H^{1}(\mathbb{R})$.这里为以后起见注意 $|H_{\varepsilon}f(x)| \leq 2C_{1}(A_{1} + A_{2})(1 + |x|)^{-2}$($\varepsilon > 0$).
\end{example}

由 $H^{1}(\mathbb{R})$ 的定义,若 $f \in H^{1}(\mathbb{R})$,则 $Hf \in L^{1}(\mathbb{R})$ 且 $-i\xi|\xi|^{-1}\mathcal{F}f(\xi) = \mathcal{F}(Hf)(\xi)$,故 $-\mathcal{F}(f)(\xi) = -i\xi|\xi|^{-1}\mathcal{F}(Hf)(\xi)$.因此,与 $f \in L^{1}(\mathbb{R})$ 合并,$Hf \in H^{1}(\mathbb{R})$ 且 $\| Hf \|_{H^{1}} = \| f \|_{H^{1}}$.总结为

\begin{theorem}\label{thm:1.12.3}
Hilbert 变换在 $H^{1}(\mathbb{R})$ 上有界,且
\begin{equation*}
\| Hf \|_{H^{1}} = \| f \|_{H^{1}} \quad (f \in H^{1}(\mathbb{R})).
\end{equation*}
\end{theorem}

下面讨论在 $H^{1}(\mathbb{R})$ 中稠密的好的空间.

\begin{lemma}\label{lem:1.12.4}
$H_{00}^{1}(\mathbb{R}) = \{f \in L^{1}(\mathbb{R});\ \hat{f} \in C^{\infty}(\mathbb{R}) \text{ 且 } \hat{f} \text{ 的支集是有界闭集且不包含原点}\}$ 在 $H^{1}(\mathbb{R})$ 中稠密.因此,$\{f \in \mathcal{S}(\mathbb{R});\ \int_{\mathbb{R}}f(x)\dd x = 0\}$ 在 $H^{1}(\mathbb{R})$ 中稠密,$H^{1}(\mathbb{R}) \cap L^{2}(\mathbb{R})$ 也在 $H^{1}(\mathbb{R})$ 中稠密.
\end{lemma}

\begin{proof}
设 $f \in H^{1}(\mathbb{R})$,任取 $\varepsilon > 0$.由 $\mathcal{F}(Hf)(\xi) = -i\xi|\xi|^{-1}\hat{f}(\xi)$,且 $Hf \in L^{1}(\mathbb{R})$,故 $-i\xi|\xi|^{-1}\hat{f}(\xi)$ 是连续函数.因此 $\hat{f}(0) = 0$,$\mathcal{F}(Hf)(0) = 0$.于是取 $\varphi \in \mathcal{S}(\mathbb{R})$ 使 $\hat{\varphi}(\xi) = 1$($|\xi| \leq 1$),$\hat{\varphi}(\xi) = 0$($|\xi| \geq 2$).首先证明 $\lim_{A \to \infty}\| \varphi_{A} * f \|_{L^{1}} = 0$.由 $\hat{f}(0) = 0$ 得 $\int_{\mathbb{R}}f(x)\dd x = 0$,故
\begin{align*}
\displaystyle\int_{\mathbb{R}}|\varphi_{A} * f(x)|\dd x & = \int_{\mathbb{R}}\left|\int_{\mathbb{R}}\left(\varphi_{A}(x - y) - \varphi_{A}(x)\right)f(y)\dd y\right|\dd x \\
& \leq \displaystyle\int_{\mathbb{R}}\left(\int_{\mathbb{R}}\frac{1}{A}\left|\varphi\left(\frac{x - y}{A}\right) - \varphi\left(\frac{x}{A}\right)\right|\dd x\right)|f(y)|\dd y \\
& = \displaystyle\int_{\mathbb{R}}\left(\int_{\mathbb{R}}|\varphi(x - y/A) - \varphi(x)|\dd x\right)|f(y)|\dd y.
\end{align*}
这里,$\int_{\mathbb{R}}|\varphi(x - y/A) - \varphi(x)|\dd x \leq 2\int_{\mathbb{R}}|\varphi(x)|\dd x$,且因 $\varphi \in L^{1}(\mathbb{R})$,$\lim_{A \to \infty}\int_{\mathbb{R}}|\varphi(x - y/A) - \varphi(x)|\dd x = 0$.因此由 Lebesgue 收敛定理,
\begin{equation*}
\lim_{A \to \infty}\int_{\mathbb{R}}\left(\int_{\mathbb{R}}|\varphi(x - y/A) - \varphi(x)|\dd x\right)|f(y)|\dd y = 0.
\end{equation*}
由此 $\lim_{A \to \infty}\| \varphi_{A} * f \|_{L^{1}} = 0$.同样,$\lim_{A \to \infty}\| \varphi_{A} * (Hf) \|_{L^{1}} = 0$.于是取 $A > 0$ 使 $\| \varphi_{A} * f \|_{L^{1}} < \varepsilon$,$\| \varphi_{A} * (Hf) \|_{L^{1}} < \varepsilon$,令 $g = f - \varphi_{A} * f$,则 $\| g - f \|_{L^{1}} < \varepsilon$,$\| Hg - Hf \|_{L^{1}} = \| \varphi_{A} * (Hf) \|_{L^{1}} < \varepsilon$.另外,由 $\hat{\varphi}(\xi) = 1$($|\xi| \leq 1$)与 $\hat{g}(\xi) = \hat{f}(\xi) - \hat{\varphi}(A\xi)\hat{f}(\xi)$,$\hat{g}(\xi) = 0$($|\xi| \leq 1/A$).其次,由逼近定理~\ref{thm:1.1.9},$\lim_{t \to 0}\| \varphi_{t} * g - g \|_{L^{1}} = 0$,$\lim_{t \to 0}\| \varphi_{t} * (Hg) - Hg \|_{L^{1}} = 0$,故取 $0 < a < A$ 使 $\| \varphi_{a} * g - g \|_{L^{1}} < \varepsilon$,$\| \varphi_{a} * (Hg) - Hg \|_{L^{1}} < \varepsilon$,令 $h = \varphi_{a} * g$.由 $\hat{h}(\xi) = \hat{\varphi}(a\xi)\hat{g}(\xi)$,$\hat{g}(\xi) = 0$($|\xi| \leq 1/A$),$\hat{\varphi}(\xi) = 0$($|\xi| \geq 2$),故 $\hat{h}(\xi) = 0$($|\xi| \leq 1/A$ 或 $|\xi| \geq 2/a$).注意 $\hat{h} \in L^{1}(\mathbb{R})$.

最后,取 $\psi \in \mathcal{S}(\mathbb{R})$ 使 $\psi(\xi) = 0$($|\xi| \geq 1$),$\mathcal{F}^{-1}\psi(0) = (2\pi)^{-1}\int_{\mathbb{R}}\psi(\xi)\dd\xi = 1$.则 $|h(x)\mathcal{F}^{-1}\psi(tx) - h(x)| \leq (\| \psi \|_{L^{1}} + 1)|h(x)|$,且对每个 $x \in \mathbb{R}$,$\lim_{t \to 0}\left(h(x)\mathcal{F}^{-1}\psi(tx) - h(x)\right) = 0$,故由 Lebesgue 收敛定理,$\lim_{t \to 0}\int_{\mathbb{R}}|h(x)\mathcal{F}^{-1}\psi(tx) - h(x)|\dd x = 0$.另外,因 $\hat{h} \in L^{1}(\mathbb{R})$,$\mathcal{F}^{-1}(\hat{h} * \psi_{t})(x) = 2\pi\mathcal{F}^{-1}\hat{h}(x)\mathcal{F}^{-1}(\psi_{t})(x) = 2\pi h(x)\mathcal{F}^{-1}\psi(tx)$,因此 $\mathcal{F}\left(2\pi h(\cdot)\mathcal{F}^{-1}\psi(t\cdot)\right) = \hat{h} * \psi_{t}$.由此 $\mathcal{F}\left(H\left(h(\cdot)\mathcal{F}^{-1}\psi(t\cdot)\right)\right)(\xi) = -i\xi|\xi|^{-1}\hat{h} * \psi_{t}(\xi)$.

现在,设 $0 < t < 1/(2A)$.若 $|\xi|, |\eta| < 1/(2A)$,则 $|\xi - \eta| < 1/A$,且 $\hat{h}(\xi - \eta) = 0$($|\xi - \eta| < 1/A$),$\psi_{t}(\eta) = 0$($|\eta| > t$),故 $\int_{\mathbb{R}}\hat{h}(\xi - \eta)\psi_{t}(\eta)\dd\eta = 0$($|\xi| < 1/(2A)$).同样,若 $|\xi| > 2/a + 1/(2A)$,$|\eta| < 1/(2A)$,则 $|\xi - \eta| > 2/a$,且 $\hat{h}(\xi - \eta) = 0$($|\xi - \eta| > 2/a$),$\psi_{t}(\eta) = 0$($|\eta| > t$),故 $\int_{\mathbb{R}}\hat{h}(\xi - \eta)\psi_{t}(\eta)\dd\eta = 0$($|\xi| > 2/a + 1/(2A)$).于是取 $k \in \mathcal{S}(\mathbb{R})$ 使 $k(\xi) = -i\xi|\xi|^{-1}$($1/(2A) < |\xi| < 2/a + 1/(2A)$),令 $K(x) = \mathcal{F}^{-1}k(x)$.则 $K \in \mathcal{S}(\mathbb{R})$,且因 $\hat{h}(\xi) = 0$($|\xi| \leq 1/A$ 或 $|\xi| \geq 2/a$),$\hat{h} * \psi_{t}(\xi) = 0$($|\xi| \leq 1/(2A)$ 或 $|\xi| \geq 2/a + 1/(2A)$),故 $k(\xi)\hat{h}(\xi) = -i\xi|\xi|^{-1}\hat{h}(\xi)$,$k(\xi)\hat{h} * \psi_{t}(\xi) = -i\xi|\xi|^{-1}\hat{h} * \psi_{t}(\xi)$,
\begin{equation*}
Hh(x) = K * h(x), \quad H\left(h(\cdot)\mathcal{F}^{-1}\psi(t\cdot)\right)(x) = K * \left(h\mathcal{F}^{-1}(\psi_{t})\right)(x).
\end{equation*}
由此
\begin{align*}
\| Hh - H\left(h(\cdot)\mathcal{F}^{-1}\psi(t\cdot)\right) \|_{L^{1}} & = \| K * \left(h - h\mathcal{F}^{-1}(\psi_{t})\right) \|_{L^{1}} \\
& \leq \| K \|_{L^{1}}\| h - h\mathcal{F}^{-1}(\psi_{t}) \|_{L^{1}}.
\end{align*}
因此,因 $\lim_{t \to 0}\| h - h\mathcal{F}^{-1}(\psi_{t}) \|_{L^{1}} = 0$,取 $0 < b < 1/(2A)$ 使
\begin{equation*}
\| h - h\mathcal{F}^{-1}(\psi_{b}) \|_{L^{1}}, \quad \| Hh - H\left(h(\cdot)\mathcal{F}^{-1}\psi(b\cdot)\right) \|_{L^{1}} < \varepsilon,
\end{equation*}
令 $f_{0} = h\mathcal{F}^{-1}(\psi_{b})$.则
\begin{equation*}
\hat{f}_{0}(\xi) = 0 \quad \left(|\xi| < 1/(2A) \text{ 或 } |\xi| > 2/a + 1/(2A)\right),
\end{equation*}
\begin{equation*}
\| f - f_{0} \|_{L^{1}} \leq \| f - g \|_{L^{1}} + \| g - h \|_{L^{1}} + \| h - f_{0} \|_{L^{1}} < 3\varepsilon,
\end{equation*}
\begin{align*}
\| Hf - Hf_{0} \|_{L^{1}} &\leq \| Hf - Hg \|_{L^{1}} + \| Hg - Hh \|_{L^{1}} + \| Hh - Hf_{0} \|_{L^{1}} \\
&< 3\varepsilon.
\end{align*}
即 $\| f - f_{0} \|_{H^{1}} < 6\varepsilon$.另外,因 $\hat{f}_{0}(\xi) = \hat{h} * \psi_{b}(\xi)$,$\psi_{b} \in \mathcal{S}(\mathbb{R})$,$\hat{h} \in L^{1}(\mathbb{R})$,故 $\hat{f}_{0} \in C^{\infty}(\mathbb{R})$.由这些可知 $H_{00}^{1}(\mathbb{R})$ 在 $H^{1}(\mathbb{R})$ 中稠密.
\end{proof}

其次,考察复平面上半平面中的正则函数与 Hardy 空间的关系.
\begin{equation*}
P(x) = \frac{1}{\pi}\frac{1}{1 + x^{2}}, \quad P(x, y) = \frac{1}{\pi}\frac{y}{y^{2} + x^{2}}, \quad Q(x, y) = \frac{1}{\pi}\frac{x}{y^{2} + x^{2}} \quad (y > 0,\ x \in \mathbb{R})
\end{equation*}
令($y > 0$,$x \in \mathbb{R}$).此时,$P_{t}(x) = P(x/t)/t = P(x, t)$,且由 $\int_{0}^{\infty}e^{-t\xi}e^{ix\xi}\dd\xi = (t - ix)^{-1}$,
\begin{equation*}
\mathcal{F}^{-1}(e^{-t|\xi|})(x) = P(x, t), \quad \mathcal{F}^{-1}\left(-i\frac{\xi}{|\xi|}e^{-t|\xi|}\right)(x) = Q(x, t) \quad (t > 0).
\end{equation*}
对 $g \in L^{1}(\mathbb{R})$,$h \in L^{2}(\mathbb{R})$,$\mathcal{F}(g * h)(\xi) = \mathcal{F}(g)(\xi)\mathcal{F}(h)(\xi)$ 容易证明,故注意 $Q(x, t)$ 关于 $x$ 属于 $L^{2}(\mathbb{R})$,则若 $f \in H^{1}(\mathbb{R})$ 且为实值函数,
\begin{equation*}
\begin{aligned}
\mathcal{F}(P_{t} * (Hf))(\xi) &= \mathcal{F}(P_{t})(\xi)\mathcal{F}(Hf)(\xi) = e^{-t|\xi|}\cdot\left(-i\dfrac{\xi}{|\xi|}\right)\hat{f}(\xi) \\
&= \mathcal{F}(Q(\cdot, t))(\xi)\hat{f}(\xi) = \mathcal{F}(Q(\cdot, t) * f)(\xi),
\end{aligned}
\end{equation*}
故对 $z = x + iy$,
\begin{equation*}
F(z) := \left(P_{y} * (f + iHf)\right)(x) = P(\cdot, y) * f(x) + iQ(\cdot, y) * f(x).
\end{equation*}
$P(x, y)$,$Q(x, y)$ 满足 Cauchy--Riemann 关系式容易验证,且注意 $f \in L^{1}(\mathbb{R})$,可知 $P(\cdot, y) * f(x)$,$Q(\cdot, y) * f(x)$ 满足 Cauchy--Riemann 关系式.因此 $F(z)$ 在上半平面正则.进而,
\begin{equation*}
\int_{\mathbb{R}}|F(x + iy)|\dd x \leq \int_{\mathbb{R}}P_{y}(x)\dd x\int_{\mathbb{R}}|f(x) + iHf(x)|\dd x \leq \| f \|_{H^{1}} \quad (y > 0).
\end{equation*}
实际上,其逆也成立.这里不给出证明,但下式成立.

\begin{theorem}\label{thm:1.12.5}
若 $F(z)$ 在上半平面正则且
\begin{equation*}
\| F \|_{H^{1}(\mathbb{R}_{+}^{2})} := \sup_{y > 0}\int_{\mathbb{R}}|F(x + iy)|\dd x < +\infty,
\end{equation*}
则 $f(x) = \lim_{y \to 0}\operatorname{Re}F(x + iy)$ 对几乎所有的 $x \in \mathbb{R}$ 存在,$f \in H^{1}(\mathbb{R})$,且
\begin{equation*}
F(z) = P_{y} * (f + iHf)(x),
\end{equation*}
$\| f \|_{H^{1}}$ 与 $\| F \|_{H^{1}(\mathbb{R}_{+}^{2})}$ 等价.
\end{theorem}

上述定义中,把 $|F(x + iy)|$ 换成 $|F(x + iy)|^{p}$($p > 0$)者,就是上半平面的一般 Hardy 空间 $H^{p}(\mathbb{R}_{+}^{2})$,它是单位圆盘上的经典 Hardy 空间 $H^{p}(U)$ 的上半平面版本.

下面不加证明地列出 $H^{1}(\mathbb{R})$ 的其他特征刻画.

\begin{theorem}\label{thm:1.12.6}
对 $f \in L^{1}(\mathbb{R})$,下面四者等价.

(i) $f \in H^{1}(\mathbb{R})$.

(ii) $\displaystyle\int_{\mathbb{R}}\sup_{t>0}|P_{t} * f(x)|\dd x < +\infty$.

(iii) $\displaystyle\int_{\mathbb{R}}\sup_{|u - x| < t}|P_{t} * f(u)|\dd x < +\infty$.

(iv)(原子分解)存在数列 $\{\lambda_{j}\}_{j=1}^{\infty}$ 与函数列 $a_{j}(x)$ 满足:

(a) $\displaystyle\sum_{j=1}^{\infty}|\lambda_{j}| < \infty$;

(b) $a_{j}(x)$ 的支集包含在某个区间 $Q_{j}$ 中,且 $\int_{Q_{j}}a_{j}(x)\dd x = 0$,$\| a_{j} \|_{L^{\infty}} \leq |Q_{j}|^{-1}$;

(c) $f(x) = \displaystyle\sum_{j=1}^{\infty}\lambda_{j}a_{j}(x)$ in $L^{1}(\mathbb{R})$.

此时,(ii),(iii) 的积分值,以及 (iv) 中 $f$ 的表示中 $\displaystyle\sum_{j=1}^{\infty}|\lambda_{j}|$ 的下确界,分别与 $\| f \|_{H^{1}}$ 等价.
\end{theorem}
